\documentclass[10pt,a4paper]{article}
\usepackage[T1]{fontenc}
\usepackage{lmodern}
\usepackage{amsmath,amssymb,amsthm,mathtools}
\usepackage[margin=24mm]{geometry}
\usepackage{microtype}
\usepackage{booktabs,array,longtable}
\usepackage{xcolor}
\usepackage{hyperref}
\usepackage{fancyhdr}
\hypersetup{colorlinks=true,linkcolor=black,citecolor=black,urlcolor=blue,pdfauthor={Matthew J. Colbrook}}
\pagestyle{fancy}
\fancyhf{}
\fancyhead[L]{\small Resolution / MI-06}
\fancyhead[R]{\small 11 September 2026}
\fancyfoot[C]{\thepage}
\setlength{\headheight}{14pt}
\setlength{\parskip}{4pt}
\setlength{\parindent}{0pt}
\newtheorem{theorem}{Theorem}[section]
\newtheorem{proposition}[theorem]{Proposition}
\newtheorem{lemma}[theorem]{Lemma}
\newtheorem{corollary}[theorem]{Corollary}
\theoremstyle{remark}
\newtheorem{remark}[theorem]{Remark}
\newcommand{\M}{\mathbb M}
\newcommand{\C}{\mathbb C}
\newcommand{\R}{\mathbb R}
\newcommand{\tr}{\operatorname{tr}}
\newcommand{\diag}{\operatorname{diag}}
\newcommand{\spec}{\operatorname{spec}}
\newcommand{\rank}{\operatorname{rank}}
\newcommand{\abs}[1]{\lvert #1\rvert}
\newcommand{\norm}[1]{\lVert #1\rVert}
\newcommand{\opnorm}[1]{\lVert #1\rVert_{\infty}}
\newcommand{\schatten}[2]{\lVert #1\rVert_{#2}}
\newcommand{\symmod}{\mathcal S}
\newcommand{\maxmod}{\mathcal M}
\newcommand{\draftnotice}{\begin{center}\small\textit{Proposed mathematical resolution. Complete proof supplied; not submitted or peer reviewed.}\end{center}}

\usepackage{xurl}
\setlength{\emergencystretch}{3em}
\allowdisplaybreaks[1]

\title{MI-06: Resolution}
\author{Matthew J. Colbrook\\
\small Department of Applied Mathematics and Theoretical Physics\\
\small University of Cambridge, Cambridge, United Kingdom\\
\small\href{mailto:m.colbrook@damtp.cam.ac.uk}{m.colbrook@damtp.cam.ac.uk}}
\date{11 September 2026}
\begin{document}
\maketitle
\textbf{Repository status: Solved.}
The complete argument received an independent Codex-agent PASS review
against the exact catalog target.
The original draft was AI-assisted; the reviewed proof below is unchanged.
This is independent agent verification, not external human peer review or
formal proof certification. \href{https://github.com/ajt60gaibb/OpenProblemsInNLA/blob/main/references/colbrook-matrix-2026-09-11/verification/reviews/MI-06-review.md}{Detailed review and scope}.
\par\medskip
% BEGIN REVIEWED PROOF
\section{MI-06: no finite two-unitary constant for the arithmetic symmetric modulus}
\label{sec:mi06}

Define the arithmetic symmetric modulus by
\[
 \symmod(X)=\frac{|X|+|X^*|}{2}.
\]
The target proposes that for every $X,Y$ there are unitaries $U,V$ with
\[
 \symmod(X+Y)\leq\sqrt2\bigl(U\symmod(X)U^*+V\symmod(Y)V^*\bigr).
\]
The following result disproves that assertion and the more general finite-constant question in \cite[Conjecture 1.6 and equation (1.4)]{Zhang2026}.

\begin{theorem}\label{thm:mi06}
Even in dimension three, there is no finite nonnegative constant $C$ such that every pair $X,Y\in\M_3(\C)$ admits unitaries $U,V$ satisfying
\begin{equation}\label{eq:mi06-target}
 \symmod(X+Y)\leq C\bigl(U\symmod(X)U^*+V\symmod(Y)V^*\bigr).
\end{equation}
The failure is witnessed by real rank-one matrices.
\end{theorem}

\begin{proof}
For $t>0$, let
\[
 X=e_1(e_1+te_2)^*
 =\begin{pmatrix}1&t&0\\0&0&0\\0&0&0\end{pmatrix},\qquad
 Y=-(e_1+te_3)e_1^*
 =\begin{pmatrix}-1&0&0\\0&0&0\\-t&0&0\end{pmatrix}.
\]
Put $s=\sqrt{1+t^2}$. If $Z=uv^*$ has rank one, then
\[
 |Z|=\frac{\|u\|}{\|v\|}vv^*,\qquad
 |Z^*|=\frac{\|v\|}{\|u\|}uu^*.
\]
Applying this identity shows that both $\symmod(X)$ and $\symmod(Y)$ have eigenvalues
\[
 a=\frac{s+1}{2},\qquad b=\frac{s-1}{2},\qquad0.
\]
On the other hand,
\[
 X+Y=t(e_1e_2^*-e_3e_1^*)
\]
has right and left moduli $t\diag(1,1,0)$ and $t\diag(1,0,1)$, respectively. Hence
\begin{equation}\label{eq:mi06-sum}
 \symmod(X+Y)=\diag(t,t/2,t/2)\geq(t/2)I_3.
\end{equation}

Fix arbitrary unitaries $U,V$. Choose a unit vector $w$ orthogonal to a top eigenvector of $U\symmod(X)U^*$ and to a top eigenvector of $V\symmod(Y)V^*$. Such a vector exists because the ambient dimension is three. The eigenvalue lists above imply
\[
 w^*\bigl(U\symmod(X)U^*+V\symmod(Y)V^*\bigr)w\leq 2b=s-1.
\]
Combining this with \eqref{eq:mi06-sum}, any domination \eqref{eq:mi06-target} would force
\[
 \frac t2\leq C(s-1),\qquad
 C\geq\frac{t}{2(\sqrt{1+t^2}-1)}
 =\frac{\sqrt{1+t^2}+1}{2t}.
\]
The final expression tends to $+\infty$ as $t\downarrow0$.
\end{proof}

\begin{corollary}[A small exact counterexample to $\sqrt2$]\label{cor:mi06-rational}
Take $t=3/4$ in the preceding construction. Then
\[
 \spec\symmod(X)=\spec\symmod(Y)=\{9/8,1/8,0\},
 \qquad \symmod(X+Y)=\diag(3/4,3/8,3/8).
\]
For arbitrary $U,V$, the quadratic form of the proposed right-hand side on the vector $w$ used above is at most $\sqrt2/4<3/8$. Thus the conjectured domination fails.
\end{corollary}

For direct verification, the two individual symmetric moduli in this example are rational matrices:
\[
 \symmod(X)=\begin{pmatrix}41/40&3/10&0\\3/10&9/40&0\\0&0&0\end{pmatrix},\qquad
 \symmod(Y)=\begin{pmatrix}41/40&0&3/10\\0&0&0\\3/10&0&9/40\end{pmatrix}.
\]
This obstruction concerns order domination by two separately rotated summands. It does not contradict a unitarily invariant norm inequality for their sum, and it does not identify the arithmetic modulus with the quadratic symmetric modulus.

% END REVIEWED PROOF
\begin{thebibliography}{1}

\bibitem{Zhang2026}
Teng Zhang.
\newblock Operator symmetric moduli and sharp triangle inequalities.
\newblock arXiv:2603.01046, 2026.
\newblock Version consulted: v1; Conjecture 1.6, equation (1.4), Problems 1.11
  and 1.17.
\newblock URL: \url{https://arxiv.org/abs/2603.01046}.

\end{thebibliography}

\end{document}
